% Generated draft: requires mathematical and source review.
\documentclass[11pt]{article}
\usepackage[margin=25mm]{geometry}
\usepackage{fontspec}
\setmainfont{Latin Modern Roman}
\usepackage{amsmath,amssymb,mathtools,unicode-math}
\setmathfont{Latin Modern Math}
\usepackage{enumitem,needspace,xurl,hyperref}
\setlength{\parindent}{0pt}
\setlength{\parskip}{5pt}
\setlength{\emergencystretch}{3em}
\newcommand{\N}{\mathbb{N}}
\newcommand{\Z}{\mathbb{Z}}
\newcommand{\Q}{\mathbb{Q}}
\newcommand{\R}{\mathbb{R}}
\newcommand{\C}{\mathbb{C}}
\newcommand{\F}{\mathbb{F}}
\newcommand{\A}{\mathbb{A}}
\newcommand{\PP}{\mathbb{P}}
\newcommand{\E}{\mathbb{E}}
\newcommand{\eps}{\varepsilon}
\newcommand{\dd}{\,\mathrm{d}}
\begin{document}
\section*{UnsolvedMath ID: 20002717}
\textbf{Source problem code:} catalogue entry
\subsection{Definitions and mathematical statement}
J. McCammond's source question calls an edge-labeled structure quasi-Garside if it has the intended Garside properties except the lattice property. For a group $G$ with a finite conjugacy-closed generating set $T$, it asks for natural conditions explaining which Cayley intervals are lattices, and whether this can be understood group-theoretically rather than through a reflection representation. The full quasi-Garside axioms and inverse/positive-word convention are not given in the extract.

For the aid's concrete cases take $T=T^{-1}$ and its word metric. The geodesic interval from $1$ to $g$ has
\[
 [1,g]_T=\{u:\ell_T(u)+\ell_T(u^{-1}g)=\ell_T(g)\},\qquad
 u\le_Tv\iff\ell_T(u)+\ell_T(u^{-1}v)=\ell_T(v).
\]
A lattice has a unique least upper bound and greatest lower bound for every pair. Left translation reduces general Cayley intervals to these based intervals.

The aid's positive criterion fixes a marked direct product
\[
 G=F\times C_1\times\cdots\times C_r,
\]
where $F$ is any finite group and each $C_i$ is finite or infinite cyclic. Use every nonidentity element of $F$ as a generator in its coordinate, and the two standard cyclic directions in each cyclic coordinate, omitting identities/duplicates. This finite symmetric marking is conjugacy-closed. Word length is the sum of coordinate lengths, with the finite coordinate contributing $0$ or $1$. Every interval is a direct product of its coordinate intervals: a point or two-chain for $F$; a chain for an infinite cyclic coordinate; and a shortest-arc chain in a finite cycle, except that an even cycle's antipodal endpoint gives two arms with common bottom and top. The two-element cyclic case reduces to a two-chain. Each factor is a lattice, so every such marked-product interval is a lattice. This is not a claim for arbitrary semidirect products or arbitrary markings of the same group.

The counterexample retains the same hypotheses on finiteness, symmetry, and conjugacy closure. In additive $\Z^2$ use all eight king moves
\[
 T=\{-1,0,1\}^2\setminus\{(0,0)\},\qquad\ell_T(a,b)=\max\{|a|,|b|\}.
\]
In $[0,(3,0)]_T$, the atoms $(1,0),(1,1)$ have two distinct rank-two common upper bounds $(2,0),(2,1)$. Both are minimal common upper bounds, so the atoms have no join; dually those rank-two elements have no meet. Thus the interval is not a lattice. The standard axial marking on $\Z^2$, covered by the positive criterion, instead gives lattice intervals, showing that the marked pair $(G,T)$ matters. This does not rule out group-structural criteria within more restricted natural marking classes.

The original remark credits Brady--Watt with a uniform lattice proof for $NC_W$ using the real-reflection realization and asks how far methods independent of that realization extend. The aid's direct-product family and bowtie obstruction do not supply a general classification of quasi-Garside interval lattices.
\subsection{Short English statement}
The supplied result classifies intervals for one explicit direct-product marking and gives a nonlattice king-move interval in $\Z^2$. It shows that finite conjugacy-closed generators alone do not force the lattice property, and that the chosen marking matters even for a fixed abelian group.
\subsection{Sources}
\begin{itemize}
\item Local UnsolvedMath registry, permanent ID 20002717, source code catalogue entry; created and updated 14 August 2026. Canonical \texttt{aim-probability-notes.json}, zero-based \texttt{notes[158]}, tag \texttt{problem}.
\item AIM ``Braid groups, clusters and free probability,'' Problem 4.3 attributed to J. McCammond, as restored by the supplied verification: \url{https://aimath.org/WWN/braidgroups/braidgroups.pdf}. The split ``Problem 4.''/``3'' is OCR damage and terminal 10 is a page number. Statement metadata is exact.
\item The original remark names Tom Brady and Colum Watt [14] for their uniform noncrossing-lattice proof, but supplies no full bibliography. The original literature field is empty. Research attempt 1: \texttt{partial\_result}, problem status at run \texttt{partially\_solved}, rigor label \texttt{proved}, novelty confidence low.
\item Ulam AI / UnsolvedMath Contributors, CC BY 4.0. Machine-generated research aid, not peer reviewed; independently verify before citation. No complete authoritative references or full quasi-Garside axioms are supplied.
\end{itemize}
\end{document}
